%% =====================================================================
%%  metadata.tex  --  DATOS del documento (lo unico que se edita aqui)
%% ---------------------------------------------------------------------
%%  El estilo lo pone el framework; este archivo solo declara datos.
%%  Los defaults del despacho (autor, cargo, lema, siglas, distrito)
%%  ya vienen del nucleo: sobrescribir solo lo que cambie.
%% =====================================================================

%% Genero del documento (cambia la palabra grande de la cabecera):
%%   Comunicado (defecto) | Pronunciamiento | Nota de Prensa | Anuncio
% \setmeta{tipodoc}{Pronunciamiento}

\setmeta{numero}{001}
\setmeta{anio}{2026}
% \setmeta{siglas}{ABC/CR}          % default del nucleo
%% Documento sin numero: comentar \setmeta{numero} (la linea desaparece)

\setmeta{lugar}{Lima}
\setmeta{fecha}{6 de agosto del 2026}   % comentar para usar la fecha de hoy

%% Nombre institucional del PDF (estandar del despacho):
%%    TIPO N°XXX-AAAA - Sumilla.pdf   (se compone al compilar)
%% Descomentar y poner la descripcion breve del tema:
% \setmeta{sumilla}{Encuentro por la educación ayacuchana}
%% O bien un nombre totalmente a medida:
% \setmeta{nombrearchivo}{COMUNICADO N°001-2026 - Encuentro educación}

%% Contacto de prensa propio (si difiere del despacho); lo usa
%% \contactoprensa en las notas de prensa:
% \setmeta{prensanombre}{Oficina de Prensa del Despacho}
% \setmeta{prensacorreo}{prensa@ejemplo.pe}
% \setmeta{prensacelular}{999 999 999}
